\documentclass[11pt]{article}
\usepackage[margin=1in]{geometry}
\usepackage{parskip}

\title{ECS272 Reading Report 6}
\author{Pablo Rodriguez}
\date{}

\begin{document}

\maketitle

\section*{Paper: More Than Telling a Story: Transforming Data into Visually Shared Stories}

\subsection*{Summary}

This paper is about how visualization storytelling is more than just making a final chart. The authors say it should be seen as a full process. They also try to narrow what counts as a visual data story so it is not just any chart on the web. The process they describe has three parts. First is exploring data, which is where people look for interesting facts or excerpts. Second is making a story, which is where those excerpts get organized into a plot with a clear message. Third is telling a story, where the plot turns into actual story material like slides, videos, or interactive visuals and then it gets delivered to an audience. The paper also says this process can loop back and does not always go in a straight line. 

Another idea in the paper is that the audience, setting, and medium really matter. Those external factors affect what gets selected, how the narrative is structured, and how it is presented. A live presentation and a web story are not the same thing, so the choices you make should change too. I am not totally sure if the authors would say this is the main takeaway, but that is what stood out to me. It felt like a call to study the whole pipeline from data to shared understanding instead of only studying the final visualization.

\subsection*{Question 1: What are the components of the storytelling process? Describe them.}

The process has three main components. Exploring data is where someone examines the dataset and pulls out excerpts or facts that might be worth sharing. Making a story is where those excerpts are turned into a plot by ordering them and connecting them so there is a clear message. Telling a story is where the plot becomes story material like annotated charts, narration, or interaction, and then it is delivered to an audience who forms the perceived story and might give feedback.

\subsection*{Question 2: Describe what kind of impact the audience and setting have on the storytelling process.}

The audience and setting affect the process at every stage because they shape what will make sense to people. The paper says the target audience is considered throughout, which means the data excerpts, plot structure, and presentation style all adjust to who is receiving the story. Setting includes time and place and also how much the audience can participate. A synchronous, colocated setting allows for back and forth interaction and changes how you pace the story. An asynchronous, distributed setting means you need clearer annotations and stronger narrative guidance since there is no live presenter. The medium also matters because a static image, a video, and an interactive system all support different kinds of storytelling.

\subsection*{Question 3: Are there any differences in the storytelling process between explaining scientific findings to novices at a museum on a touch table versus creating a data story on the New York Times webpage. Why or why not?}

Yes, I think there are differences because the setting and medium are really different. A museum touch table is synchronous and colocated with high participation, so the presenter can see reactions, answer questions, and adjust the story in real time. That pushes the story material to be more interactive and exploratory with guidance built into the interface. A New York Times webpage is asynchronous and distributed with low participation, so the reader is alone and the story needs stronger annotation and tighter sequencing. The overall steps of exploring data, making a story, and telling a story are still there, but the constraints and design choices change because the audience and setting are different.

\end{document}
